\chapter{Spin Geometry}

\dropcap{T}{he purpose} of this chapter is to introduce the algebraic and geometric objects that appear in the Seiberg-Witten equations. Clifford algebras, spin groups, spinor bundles, and Dirac operators provide the language in which the equations are formulated.

\section{Clifford algebras}
\subsection{Definition and basic properties}
\begin{definition}[Clifford algebra]
Let $(V,q)$ be a quadratic vector space. The Clifford algebra $\Cl(V,q)$ is the quotient of the tensor algebra $T(V)$ by the ideal generated by elements of the form
\[
  v\otimes v + q(v)1.
\]
\end{definition}

\begin{notation}
We denote by $\R^{r,s}$ the semi-Riemannian vector space $\R^{r+s}$ with $r$ positive and $s$ negative directions. We write $\Cl(r,s)$ for the corresponding Clifford algebra.
\end{notation}

\begin{example}[$\Cl(1)$ and $\Cl(2)$]
For $n=1$, the real Clifford algebra generated by one element with relation $e^2=-1$ is isomorphic to $\C$. For $n=2$, one obtains an algebra isomorphic to the quaternions $\mathbb{H}$.
\end{example}

\subsection{Complex(ified) Clifford algebras}
Complex Clifford algebras are often simpler to classify and are central for describing spin representations.

\section{The Spin group}
\subsection{Definition and examples}
\begin{definition}[Spin group]
The group $\Spin(n)$ is the subgroup of the Clifford algebra generated by products of unit vectors and fitting into the double cover
\[
1\to \Ztwo\to \Spin(n)\to \operatorname{SO}(n)\to 1.
\]
\end{definition}

\subsection{When $n=4$}
Dimension four is special because $\Spin(4)$ decomposes in a way that makes four-dimensional gauge theory particularly rich.

\subsection{Lie algebra structures}
Lie algebra computations identify infinitesimal versions of the spin actions used later.

\section{Representations}
\subsection{Classification and examples}
\begin{theorem}[Irreducible representations of matrix algebras]
Let $A$ be a finite direct sum of matrix algebras over a field. The irreducible representations of $A$ are the projections onto the simple summands, followed by the standard representation.
\end{theorem}

\subsection{The Spin representation}
Spin representations arise by restricting Clifford algebra representations to the spin group.

\subsection{Clifford multiplication and infinitesimal actions}
Clifford multiplication connects algebraic representations with first-order differential operators.

\section{Spin structures}
\subsection{From local to global}
\begin{definition}[Spin structure]
Let $M$ be an oriented Riemannian manifold. A spin structure is a principal $\Spin(n)$-bundle lifting the oriented orthonormal frame bundle of $M$.
\end{definition}

\subsection{Spinor bundles}
A spin structure allows one to construct spinor bundles, whose sections are spinor fields.

\section{The Dirac operator}
\subsection{Spin connections}
Connections on spinor bundles are induced by Levi-Civita connections and auxiliary choices.

\subsection{The Dirac operator}
\begin{definition}[Dirac operator]
Given a spinor bundle $S\to M$ with connection $\nabla$, the Dirac operator is locally given by
\[
  \Dslash\psi = \sum_i e_i\cdot \nabla_{e_i}\psi,
\]
where $\{e_i\}$ is a local orthonormal frame.
\end{definition}

\section{The world of $\Spinc$}
\subsection{The $\Spinc$ group}
\begin{definition}[$\Spinc$ group]
The group $\Spinc(n)$ is defined by
\[
  \Spinc(n)=\frac{\Spin(n)\times U(1)}{\{\pm(1,1)\}}.
\]
\end{definition}

\subsection{Going global: $\Spinc$-structures and spinors}
$\Spinc$ structures are weaker than spin structures and exist on a broader class of four-manifolds.

\subsection{The complex spin connection and the coupled Dirac operator}
The coupled Dirac operator is the analytic object appearing in the Seiberg-Witten equations.

\section{Final ingredient: The squaring map}
\subsection{The linear squaring map}
The squaring map sends spinors to self-dual two-forms in a way compatible with the spin representation.

\subsection{The global squaring map}
Globally, this construction produces the nonlinear term in the Seiberg-Witten equations.
